Many scientific modeling problems require forecasting stochastic dynamics from snapshot data. Here, snapshot data represent observations taken at different time points, without access to individual trajectories. And forecasting involves predicting future states beyond the observed times. For example, single-cell RNA sequencing (scRNA-seq) is widely used to study dynamic processes such as development, immune activation, and cancer progression. A scRNA-seq measurement destroys the cell, so scientists observe independent biological replicates at discrete times rather than a single replicate across multiple times. Despite the absence of individual cell trajectories, researchers often aim to forecast future cellular states; for instance, researchers are interested in forecasting differentiation outcomes of stem cells, immune cell activity after initial signal stimulation, or long-term cancer cell response to drugs with transcriptomic snapshots taken shortly after treatments. A common additional challenge is incomplete state measurement. For instance, although protein expression level mediates the dynamics of gene expression, protein level cannot be measured by scRNA-seq. 

Recent work has addressed the problem of interpolating between snapshots through Schrödinger bridge (SB) methods \citep{pavon2021data,de2021diffusion,Vargas2021,koshizuka2022neural,Wang2023} and their multi-marginal extensions \citep{shen2024multi,zhang2024joint,guan2024identifying,chen2024deep,Lavenant2021}. These methods reconstruct likely trajectories between snapshots at consecutive times. %SB methods interpolate via regularized couplings, often chosen to be entropy-regularized.
In settings where the goal is to fill in missing timepoints between observed snapshots, these techniques can perform well. However, to the best of our knowledge, no previous papers before the present work presented methodology for forecasting, ran experiments on forecasting, or had code that could immediately run for the forecasting task.

Nonetheless, we can imagine some natural extensions of existing work to forecasting. One option is to follow the SB reference dynamic beyond observed time points, but since the reference dynamic is typically chosen without considering data, we expect poor extrapolation performance. \citet{chen2024deep} and \citet{shen2024multi} incorporate momentum or a learned reference, respectively, so one might hope they could extend well to extrapolation.
%Some extensions \citep[e.g.,][]{chen2024deep} incorporate velocity or momentum to improve extrapolation. %However, most SB-based forecasting setups still rely on a fixed reference process and do not support data-adaptive learning of nonlinear, state-dependent dynamics. 
But we expect the extrapolation ability of the latter methods to be limited by an issue shared by all existing SB methods.
%by fixed reference dynamics or restrictive optimization objectives
%Our approach offers an alternative by learning a forward-simulatable SDE directly from marginal observations that enables generalization beyond the training time window. Furthermore, all SB-based methods
Namely, existing methods optimize a Kullback--Leibler divergence between data and a nominal diffusion model, and this optimization formulation requires a known volatility, fixed across states. In practice, volatility is often unknown or state-dependent, as in scRNA-seq, where noise arises from a mix of biological and technical sources. We see in our experiments that missing the state dependence or choosing an inappropriate (though standard) volatility value can lead to poor forecasts. In addition, most approaches (with the exception of \citet{shen2024multi}) do not take advantage of partial mechanistic knowledge (e.g., bounded production–decay kinetics in a biological system or vortex structure in ocean currents), even though such knowledge is often available and can meaningfully constrain the dynamics. Finally, we note that SB methods often rely on iterative Sinkhorn solvers
that offer %limited interpretability and 
few tools for diagnosing model fit. We describe additional related work in \cref{app:related-work}.

In this work, we propose a new framework, SnapMMD, that shifts the modeling focus from interpolation to accurate, interpretable forecasting. Our approach begins with the observation that in typical trajectory-inference settings \citep{Lavenant2021}, each sample can be viewed as an i.i.d.\ draw from the \emph{joint} distribution of a system's state (e.g., mRNA expression level) and the time of measurement. Then, we characterize a parametric family of stochastic differential equations (SDEs) and seek the member of this family whose implied joint distribution over state and time best matches the empirical joint distribution observed in the data. We perform this matching using maximum mean discrepancy (MMD) with a specific kernel choice, which we show enjoys a number of conceptual and computational benefits.
%, a kernel-based measure of distance between distributions that allows us to compare the model-predicted state-time distribution to the observed one without requiring access to individual trajectories or likelihoods. With a specific kernel choice, we show that this matching problem reduces to a distributional least-squares objective.
Our framework allows data to guide extrapolation beyond observed times, allows us to learn volatility rather than fixing it in advance, facilitates the use of domain knowledge when available, and yields interpretable model diagnostics, including an explicit velocity field and an $R^2$-like metric that quantifies model fit. Our framework offers the added benefit of enabling robust interpolation and forecasting even with incomplete state measurements, as in the protein expression example above.
%while still permitting data-driven inference.
%Finally, our formulation yields three key benefits: (1) accurate forecasts beyond the observed time horizon, (2) robust interpolation and model learning even with incomplete state measurements, and (3) interpretable model outputs, including an explicit velocity field and an $R^2$-like metric that quantifies model fit.
%
%This structure helps avoid some of the identifiability issues that arise in fully unstructured neural SDEs.
%
We evaluate SnapMMD across a range of synthetic and real-world systems. %, including a synthetic gene regulatory network and real time-course scRNA-seq data. %The experiments span settings with different amounts of prior knowledge. %: Lotka–Volterra the synthetic gene regulatory network have known mechanistic forms, while scRNA-seq and ocean dynamics are only partly specified. SnapMMD can incorporate whichever level of structure is available.
We find that SnapMMD consistently outperforms Schrödinger bridge baselines in all forecasting tasks --- while still
%e also assess SnapMMD on interpolation between observed snapshots;
matching or exceeding interpolation performance of both SB-based and flow-matching baselines in almost all cases. %Crucially, SnapMMD handles incomplete state measurements with ease, reconstructs an interpretable velocity field that supports downstream scientific analysis, and provides an $R^2$-style metric to evaluate model fit. 
